\section{训练神经网络}

\subsection{损失优化的目标}

训练神经网络的目标是找到使经验损失 $J(\mathbf{W})$ 最小化的权重 $\mathbf{W}^*$：

$$\mathbf{W}^* = \arg\min_{\mathbf{W}} J(\mathbf{W}) = \arg\min_{\mathbf{W}} \frac{1}{n}\sum_{i=1}^{n}\mathcal{L}\big(f(x^{(i)}; \mathbf{W}), y^{(i)}\big)$$

需要注意的是，$\mathbf{W}$ 包含网络中\textbf{所有层}的所有权重和偏置：$\mathbf{W} = \{\mathbf{W}^{(0)}, \mathbf{W}^{(1)}, \ldots\}$。

\begin{figure}[H]
\centering
\includegraphics[width=0.85\textwidth]{slides-images/slide-057.jpg}
\caption{损失函数的 landscape：目标是找到最低点\protect\footnotemark}
\end{figure}
\footnotetext{Slides 第 57 页。}

\subsection{梯度下降}

\textbf{梯度下降}（Gradient Descent）是优化损失函数的核心算法。其基本思想简洁而优雅：

\begin{importantbox}{梯度下降算法}
\begin{enumerate}
\item \textbf{随机初始化}权重 $\mathbf{W}$
\item \textbf{计算梯度}：$\frac{\partial J(\mathbf{W})}{\partial \mathbf{W}}$，梯度指向损失增大最快的方向
\item \textbf{更新权重}：沿梯度的\textbf{反方向}迈出一小步
$$\mathbf{W} \leftarrow \mathbf{W} - \eta \frac{\partial J(\mathbf{W})}{\partial \mathbf{W}}$$
\item 重复步骤 2--3 直到收敛
\end{enumerate}
其中 $\eta$ 是\textbf{学习率}（learning rate），控制每步更新的幅度。
\end{importantbox}

\begin{figure}[H]
\centering
\includegraphics[width=0.85\textwidth]{slides-images/slide-060.jpg}
\caption{梯度下降的可视化：沿梯度反方向迭代下降\protect\footnotemark}
\end{figure}
\footnotetext{Slides 第 60 页。}

梯度的直觉解释：想象你站在一个丘陵地带，蒙着眼睛，想要走到最低点。你能做的就是感受脚下地面的坡度（梯度），然后向最陡下降的方向迈一步。重复这个过程，最终你会到达某个低洼处。

\subsection{反向传播}

那么，如何高效地计算梯度 $\frac{\partial J}{\partial \mathbf{W}}$ 呢？答案是\textbf{反向传播}（Backpropagation）。

反向传播利用\textbf{链式法则}（Chain Rule），从输出层开始，逐层向回传播梯度。对于网络中的每个权重 $w$，我们可以通过链式法则将最终损失对该权重的梯度分解为若干中间梯度的乘积。

以一个简单的两层网络为例，损失 $J$ 对第一层权重 $w_1$ 的梯度可以表示为：

$$\frac{\partial J}{\partial w_1} = \frac{\partial J}{\partial \hat{y}} \cdot \frac{\partial \hat{y}}{\partial z_2} \cdot \frac{\partial z_2}{\partial z_1} \cdot \frac{\partial z_1}{\partial w_1}$$

\begin{importantbox}{反向传播的核心}
反向传播的本质就是\textbf{链式法则的递归应用}。从输出端开始计算局部梯度，然后逐层向后传播，将所有局部梯度相乘得到最终梯度。这使得在含有数百万参数的深度网络中高效计算梯度成为可能。
\end{importantbox}

现代深度学习框架（TensorFlow、PyTorch）通过\textbf{自动微分}（Automatic Differentiation）自动完成反向传播的梯度计算，开发者无需手动推导。

\begin{lstlisting}[caption={TensorFlow 和 PyTorch 中的自动微分}]
# TensorFlow: GradientTape
import tensorflow as tf
with tf.GradientTape() as tape:
    loss = compute_loss(model, x, y)
grads = tape.gradient(loss, model.trainable_variables)

# PyTorch: autograd
import torch
loss = compute_loss(model, x, y)
loss.backward()
# Gradients stored in param.grad
\end{lstlisting}

\subsection{学习率的重要性}

学习率 $\eta$ 是训练神经网络中最关键的超参数之一。它的设置直接影响训练的成败：

\begin{table}[H]
\centering
\begin{tabular}{lp{9cm}}
\toprule
\textbf{学习率} & \textbf{效果} \\
\midrule
太小 & 收敛极其缓慢，可能陷入局部最小值 \\
太大 & 训练不稳定，损失值在最小值附近反复震荡甚至发散 \\
适中 & 稳定收敛到良好的最小值 \\
\bottomrule
\end{tabular}
\caption{学习率设置对训练的影响}
\end{table}

\begin{warningbox}{学习率调节常见陷阱}
\begin{itemize}
\item 学习率过大不仅会导致震荡，还可能使损失值越来越大（发散），因为每步更新都跨越了最小值
\item 学习率过小虽然安全，但可能需要极长的训练时间，且容易被困在浅层的局部最小值中
\item 实践中常用\textbf{自适应学习率}方法（如 Adam、RMSProp），它们会根据梯度的历史信息自动调整每个参数的学习率
\end{itemize}
\end{warningbox}

\begin{importantbox}{自适应学习率算法}
实际训练中很少使用固定学习率。常用的自适应优化器包括：
\begin{itemize}
\item \textbf{SGD with Momentum}：利用梯度的指数移动平均来加速收敛
\item \textbf{Adam}（Adaptive Moment Estimation）：结合一阶矩和二阶矩估计，为每个参数独立调整学习率
\item \textbf{RMSProp}：通过梯度平方的移动平均来归一化学习率
\end{itemize}
这些算法的核心思想是：不再对所有参数使用同一个固定学习率，而是根据每个参数梯度的历史动态调整。
\end{importantbox}

\subsection{Mini-batch 梯度下降}

标准梯度下降需要在\textbf{整个数据集}上计算损失和梯度后才能进行一次权重更新，这在大规模数据集上计算量巨大。

\textbf{随机梯度下降}（Stochastic Gradient Descent, SGD）是另一个极端：每次只用\textbf{一个样本}来计算梯度和更新权重，虽然快但噪声很大。

实践中的折衷方案是 \textbf{Mini-batch 梯度下降}：每次从数据集中随机抽取一个小批量（batch）的样本来计算梯度：

$$\frac{\partial J(\mathbf{W})}{\partial \mathbf{W}} \approx \frac{1}{B}\sum_{k=1}^{B}\frac{\partial \mathcal{L}(f(x^{(k)}; \mathbf{W}), y^{(k)})}{\partial \mathbf{W}}$$

其中 $B$ 是 batch size。

\begin{table}[H]
\centering
\begin{tabular}{lp{4.5cm}p{4.5cm}}
\toprule
\textbf{方法} & \textbf{优点} & \textbf{缺点} \\
\midrule
Full Batch GD & 梯度精确，收敛稳定 & 计算开销大，内存需求高 \\
SGD ($B=1$) & 极快的更新频率 & 梯度噪声大，收敛不稳定 \\
Mini-batch SGD & 兼顾速度与稳定性，可并行化 & 需要调节 batch size \\
\bottomrule
\end{tabular}
\caption{三种梯度下降策略对比}
\end{table}

\begin{knowledgebox}{Mini-batch 的额外好处}
Mini-batch 梯度下降不仅是计算效率的折衷，其引入的梯度噪声实际上还有助于模型跳出尖锐的局部最小值，倾向于收敛到更平坦的最小值区域，而平坦的最小值通常具有更好的泛化性能。此外，GPU 硬件天然适合对 batch 数据进行并行计算，使得 mini-batch SGD 的实际训练速度远超逐样本更新。
\end{knowledgebox}

\subsection{梯度下降的挑战：局部最小值与鞍点}

在高维参数空间中，损失函数的 landscape 非常复杂。梯度下降面临几个关键挑战：

\begin{enumerate}
\item \textbf{局部最小值}（Local Minima）：损失函数可能有许多局部最小值，梯度下降可能收敛到一个不是全局最优的位置
\item \textbf{鞍点}（Saddle Points）：在高维空间中，鞍点比局部最小值更为常见。在鞍点处，梯度在某些方向为零，导致优化停滞
\item \textbf{平坦区域}（Plateaus）：损失函数的梯度接近零的区域，导致更新极其缓慢
\end{enumerate}

\begin{knowledgebox}{高维空间中的优化直觉}
在低维空间（如二维）中，局部最小值看起来很常见。但在神经网络常见的数百万维参数空间中，要形成一个真正的局部最小值需要\textbf{所有维度}同时向上弯曲，这在统计上极不可能。实际上，高维优化中更常见的问题是鞍点——某些维度向上、某些维度向下。这也解释了为什么深度学习在实践中能够成功：虽然损失 landscape 看起来复杂，但大多数"坑"都有逃逸路径。
\end{knowledgebox}

\subsection{本章小结}

训练神经网络的核心是通过梯度下降最小化损失函数。反向传播利用链式法则高效计算梯度。学习率是最关键的超参数，过大过小都会导致问题，自适应学习率算法（如 Adam）可以自动调节。Mini-batch SGD 在计算效率和梯度质量之间取得了良好平衡。高维空间中的优化虽然面临局部最小值和鞍点等挑战，但现代优化算法在实践中表现出色。

